\cotefile{82}
\scanpage{32}
encore que $E'$ opère \uline{\unsure{fidèlt|fidèlmt}} sur $E$ par $\rho_E\deleted{\unsure{'}}$, quand \deleted{\unsure{on}} $S=\operatorname{Spec}k$, $k$ corps alg. clos. Or dans ce cas on peut trouver une base $e_1,e_2,e_3$ p.r. à laquelle $q$ se met sous la forme $xy+z^2$, $\omega=e_1\wedge e_2\wedge e_3$ (quitte \inserted{au besoin} à changer $e_3$ en $-e_3$\ldots) et il suffit de prendre alors les calculs de (72)
\[
\rho_E(xX+hH+yY):
\left\{\begin{aligned}
\widetilde X&\longmapsto h\widetilde X-y\widetilde H\\
\widetilde Y&\unsure{\vdash|\longmapsto} x\widetilde H\unsure{=|-|+}h\widetilde Y\\
\widetilde H&\longmapsto-2x\widetilde X+2y\widetilde Y\deleted{\unsure{+}}
\end{aligned}\right.
\]
\[
\rho_E(xX+hH+yY)=0\quad\Longrightarrow x=y=h=0\qquad\text{o.k.}
\]

Il faut encore définir l'isom.
\[
 E\xrightarrow{\sim}\unsure{\mathfrak g|\mathfrak g'|\widetilde{\mathfrak g}}
\]
et prouver ses propriétés. Pour ceci, il faut redéfinir ou préciser, en termes élémentaires, la construction de $\widetilde G$ sur $G$ -- il nous faut pour ceci ouvrir une parenthèse
% archive sheet number: 31
